\section{What a bootstrap must reproduce}
\label{sec:bootstrap-consistency}

Separating fitted means from future payments is necessary, but it still leaves
the statistical target to be chosen. An insurer knows an observed claims
triangle. A method assessment may instead hold only its latest diagonal fixed,
allowing different earlier histories to lead to those same cumulative amounts.
These are different conditioning questions. Let $D$ denote the observed
triangle, $Q$ its latest diagonal, $R$ the future reserve payment and
$\widehat R(D)$ the estimate calculated from the triangle. A distribution of
$R-\widehat R(D)$ conditional on $Q$ allows the earlier entries of $D$, and
therefore the estimate, to vary. Conditional on the whole $D$, that particular
estimate is fixed.

A small authored example makes the distinction concrete. Suppose two earlier
histories are equally likely given the same $Q$ and produce estimates of 90
and 110. Let future payments be 80 or 120 with equal probability, independent
of those histories conditional on $Q$. The prediction error given $Q$ is then
$-30,-10,10,30$, each with probability $1/4$: its mean is zero and variance
500. Fixing the history whose estimate is 90 instead gives errors $-10,30$,
mean 10 and variance 400. Both calculations are correct for their stated
conditioning information. They cannot serve as interchangeable validation
targets. Nor does holding an estimate fixed imply that all uncertainty about
model parameters has disappeared; a Bayesian parameter distribution, for
example, would be a further explicitly specified model.

Steinmetz and Jentsch analyze a Mack bootstrap against the prediction error
conditional on the latest diagonal. Their asymptotic construction adds older
accident years and additional development observations while retaining the
current calendar boundary. It does not obtain more information merely by
increasing the number of bootstrap draws. The framework assumes positive
initial cumulative claims with finite moments, specified independence,
development factors tending to one, sufficiently fast decay of development
variances, and further moment and regularity conditions. It is distinct from
the incremental overdispersed-Poisson procedure in the preceding section
\citep[Sections 2--4, especially Assumptions 2.2--2.5]{steinmetz2024bootstrap}.

Their analysis separates an unscaled process term from a smaller estimation
term. Under the stated assumptions, the process term has a nondegenerate,
generally non-Gaussian limit. Correctly reproducing its distribution requires
the correct conditional family for future development factors; matching two
moments within a wrongly chosen family is insufficient. The original
bootstrap can reproduce that limit while missing the conditional mean and
covariance of the estimation term at its own, square-root sample-size scale
\citep[Assumptions 4.6--4.9, Theorems 4.8--4.11 and Remark 4.12]{steinmetz2024bootstrap}.

There is no contradiction in these two findings. For illustration, let a
process term be an equally likely $-1$ or $1$, and an independent estimation
term an equally likely $-1/n$ or $1/n$. Their sum has variance $1+1/n^2$.
Dropping the estimation term produces the same limiting distribution for the
unscaled sum, yet multiplying that omitted term by $n$ would have left variance
one. This example explains the distinction between scales; it is not a
derivation of the paper's Mack limit or evidence that a finite insurer triangle
is adequately approximated by it.

\subsection{Reconstructing the past requires a backward law}

The paper's proposed repair holds $Q$ fixed, simulates earlier cumulative
claims backwards, refits the resulting triangle, and centers simulated future
reserve payments at that refitted estimate. Future development uses the
originally estimated forward factors. The purpose of reconstructing the past
is to retain the effect of conditioning on the observed diagonal when
measuring estimation uncertainty. Theorem 6.3 requires Assumption 6.2,
including conditional limit properties for the backward reconstruction.
Appendix C invokes those properties; it does not establish them for every
positive distribution chosen by a programmer
\citep[Section 5, Assumption 6.2 and Appendix C]{steinmetz2024bootstrap}.

Forward conditional moments do not determine backward conditional moments.
Let $X$ be an earlier cumulative claim and $Y$ the next cumulative claim.
Suppose $X=1$ or $2$ with equal probability. Given $X=1$, let $Y=1$ or $3$,
equally likely. Given $X=2$, let $Y=3$ or $6$ with probabilities $2/3$ and
$1/3$. These positive claims satisfy
$\E[Y\mid X]=2X$ and $\operatorname{Var}(Y\mid X)=X$.
Thus the forward factor $Y/X$ has conditional mean 2 and variance $1/X$,
the form of a one-step Mack moment relation.

Now observe $Y=3$. The joint masses for $X=1$ and $X=2$ are $1/4$ and
$1/3$, respectively. Their sum is $7/12$, so Bayes' rule gives
$\Pr(X=1\mid Y=3)=3/7$ and $\Pr(X=2\mid Y=3)=4/7$.
The backward conditional mean is $11/7$, and its variance is $12/49$.
Dividing the observed cumulative amount by the forward mean factor would
instead give $3/2$. Choosing a backward factor with mean $1/2$ and variance
$1/3$ would imply a variance of 3 for $X$, also wrong for this joint law.
The example establishes the missing implication; it does not assert that this
finite distribution satisfies every assumption of an infinite-triangle theorem.

In the published paper, Section 5.1 proposes reciprocal-factor backward
moments, while Remark 5.1 and Assumption 6.2 separately require adequate
backward conditional behavior. The proposal therefore keeps such a generator
as a hypothesis requiring its own justification. Given a joint model, the
backward distribution must be derived from that model or approximated with a
checked error criterion. Choosing a gamma family and matching reciprocal
moments alone cannot establish the required conditional limit properties
\citep[Section 5.1, Remark 5.1 and Assumption 6.2]{steinmetz2024bootstrap}.

\subsection{The center and the comparison must use the same quantity}

Reserve payments exclude claims already paid. Write $q_i$ for the latest
cumulative payment in origin period $i$, $U_i^+$ for its simulated ultimate
payment, and $\widehat U_i^+$ for the ultimate payment predicted from a
reconstructed past. The reserve error must subtract like quantities:
\[
 R^+-\widehat R^+
 =\sum_i(U_i^+-q_i)-\sum_i(\widehat U_i^+-q_i)
 =\sum_i(U_i^+-\widehat U_i^+).
\]
For $q=100$, simulated ultimate payment 130 and fitted ultimate payment 120,
the reserve is 30, its estimate is 20, and their difference is 10.
Subtracting the ultimate estimate 120 directly from the reserve 30 gives
$-90$, shifting the error by the already-paid amount.

This distinction also resolves a literal inconsistency in the retained
published text. Equation (5.3), printed on page 96, writes the centering
quantity as a sum of projected ultimates without subtracting the latest
diagonal, although the preceding step defines reserves net of that diagonal.
Read literally, those two formulas produce the wrong reserve error. The
adjacent product-difference formula and decomposition (6.1) instead cancel
the diagonal and agree with the derivation above. The proposed implementation
uses that explicitly derived reserve center; it does not copy (5.3) literally.
No claim about the authors' unavailable execution code follows
\citep[Section 5.1, equation 5.3 and equation 6.1]{steinmetz2024bootstrap}.

The final journal article also changes the preprint's simulation comparison.
Section 7.1.1 generates whole triangles to assess unconditional prediction
interval coverage. Section 7.1.2 assesses conditional coverage using an
approximate backward generator in equation (7.4), based on the same
reciprocal-moment idea. The article acknowledges that the exact backward
conditional law is generally difficult to obtain. Its conditional simulation
therefore evaluates that constructed experiment; it does not by itself verify
coverage under the exact backward law implied by every forward model.
The published numerical differences remain reported simulation observations,
without a local replication or statistically supported ranking here
\citep[Section 7.1, Table 3 and Figures 1--3]{steinmetz2024bootstrap}.

For AccDSL, a reserving-method specification must identify the payment
horizon, conditioning information, process family, parameter-estimation
procedure, backward law when one is needed, and reserve-centering convention.
Validation must distinguish whole-triangle repeated-sample coverage,
coverage conditional on a diagonal, and prediction conditional on a full
observed history. A simulator with a convenient reverse rule is a candidate
model, not an independent reference for that same rule. Exact centering and
support checks precede simulations; target-specific comparison, uncertainty
estimation and downstream financial-statement checks precede any promotion.
The small triangle remains an accounting and arithmetic illustration, with
its unresolved dispersion preserved.
